% Part: methods
% Chapter: proofs
% Section: introduction

\documentclass[../../../include/open-logic-section]{subfiles}

\begin{document}

\olfileid{mth}{prf}{int}

\olsection{పరిచయం}

ప్రాథమిక తర్కంలో మీ అనుభవం వల్ల మీకు \tetoken{వ్యుత్పత్తి}{!!a{derivation}} వ్యవస్థ పరిచయమై ఉండవచ్చు---బహుశా సహజ నిగమన లేదా ఫిచ్ శైలి \tetoken{వ్యుత్పత్తి}{!!{derivation}} వ్యవస్థ, లేక నిరూపణ-వృక్ష వ్యవస్థ కావచ్చు. ఈ వ్యవస్థలలో \tetoken{ఒక సూత్రాన్ని}{!!a{formula}} నిరూపించడం లేదా ఇచ్చిన వాదన చెల్లుబాటవుతుందని చూపడం మీరు చేసి ఉంటారు. అందుకు కావలసిన చివరి ఫలితం వచ్చే వరకు వ్యవస్థ నియమాలను వర్తింపజేశారు. తర్కం \emph{గురించి} తర్కించేటప్పుడు కూడా మనం విషయాలను నిరూపిస్తాం; అయితే చాలా సందర్భాల్లో \tetoken{వ్యుత్పత్తి}{!!a{derivation}} వ్యవస్థను ఉపయోగించం. నిజానికి మనం పరిశీలించే నిరూపణలలో చాలా వాటిని పూర్తిగా మొదటిస్థాయి తర్క భాషలో కాక, కొంత సంకేత భాషతో కూడిన సహజ భాషలో (మూల రచనలో ఆంగ్లం; ఇక్కడ తెలుగు) రాస్తాం. అటువంటి నిరూపణను నిర్మించేటప్పుడు మొదట మీకు సందేహం రావచ్చు---\tetoken{వ్యుత్పత్తి}{!!a{derivation}} వ్యవస్థ లేకుండా దేనినైనా ఎలా నిరూపించాలి? ఎక్కడ మొదలుపెట్టాలి? నా నిరూపణ సరైనదో ఎలా తెలుసుకోవాలి?

నిరూపణకు ప్రయత్నించే ముందు నిరూపణ అంటే ఏమిటి, దాన్ని ఎలా నిర్మించాలి అనేవి తెలుసుకోవాలి. పేరే సూచించినట్లుగా, ఒక \emph{నిరూపణ} ఏదో విషయం సత్యమని చూపడానికి ఉద్దేశించబడింది. దీన్ని ఒక సంభాషణగా ఊహించవచ్చు---ఎవరైనా ఒక విషయం నిజమా అని అడుగుతారు; ఉదాహరణకు రెండు తప్ప మిగతా ప్రతి ప్రధాన సంఖ్య బేసిసంఖ్యేనా అని. ``అవును'' అని చెప్పడం సరిపోదు; అది \emph{ఎందుకు} నిజమో వారు అడగవచ్చు. అప్పుడు మీరు నిరూపణ ఇస్తారు.

రోజువారీ సంభాషణలో సమాధానాన్ని సూచించడం లేదా అసంపూర్ణ సమాధానం ఇవ్వడం సరిపోవచ్చు. కానీ తర్కంలోనూ గణితంలోనూ మనకు కట్టుదిట్టమైన నిరూపణ కావాలి---ఒక విషయం \emph{ఏ} సందేహానికీ తావులేకుండా సత్యమని చూపాలి. అంటే నిరూపణలోని ప్రతి దశకు సమర్థన ఉండాలి; ఆ సమర్థన తగినదై ఉండాలి (ఉదాహరణకు, వాడుతున్న పరికల్పన నిజంగా మీరు నిరూపిస్తున్న సిద్ధాంతపు ప్రకటనలో ఉండాలి; నిర్వచనాలను సరిగ్గా వర్తింపజేయాలి; ఆధారపడే నిగమనాలు సరైనవై ఉండాలి).

సాధారణంగా మనం ఒక ప్రకటనను నిరూపిస్తాం. నిరూపించే ప్రకటనలకు ప్రతిపాదనలు, సిద్ధాంతాలు, ఉపప్రమేయాలు, అనుబంధ ఫలితాలు వంటి పేర్లు పెడతాం. ప్రతిపాదన అనేది నిరూపించి నమోదు చేయదగిన ప్రాథమిక ప్రకటన: నమోదు చేసేంత ముఖ్యమైనదే, కానీ తప్పనిసరిగా లోతైనదో తరచుగా ఉపయోగించేదో కాకపోవచ్చు. సిద్ధాంతం మరింత ప్రాముఖ్యమైన ప్రతిపాదన. దాని నిరూపణను తరచుగా అనేక దశలుగా విడదీస్తారు; కొన్నిసార్లు దాన్ని మొదట నిరూపించిన వ్యక్తి పేరుతో (ఉదా., కాంటర్ సిద్ధాంతం, లొవెన్‌హైమ్--స్కోలెమ్ సిద్ధాంతం), లేదా అది తెలిపే విషయం పేరుతో (ఉదా., సంపూర్ణత సిద్ధాంతం) పిలుస్తారు. ఉపప్రమేయం మరింత ముఖ్యమైన ఫలిత నిరూపణలో ఉపయోగపడే ప్రతిపాదన లేదా సిద్ధాంతం. గందరగోళంగా అనిపించినా, కొన్ని ఉపప్రమేయాలు స్వయంగా ముఖ్యమైన ఫలితాలు; వాటిని ప్రవేశపెట్టిన వ్యక్తి పేరు కూడా వాటికి పెడతారు (ఉదా., జోర్న్ ఉపప్రమేయం). అనుబంధ ఫలితం మరొక ఫలితం నుంచి సులభంగా వచ్చేది.

నిరూపించాల్సిన ప్రకటనలో మనం ఏ రకమైన వస్తువుల గురించి నిరూపిస్తున్నామో స్పష్టం చేసే పరికల్పనలు తరచుగా ఉంటాయి. అది ``$!A$ అనేది $!B \lif !C$ రూపంలోని \tetoken{ఒక సూత్రం}{!!a{formula}} అనుకుందాం'' లేదా ``$\Gamma \Proves !A$ అనుకుందాం'' అని మొదలవవచ్చు. ఇవి ఆ ప్రతిపాదన, సిద్ధాంతం లేదా ఉపప్రమేయపు \emph{పరికల్పనలు}; నిరూపణలో అవి సత్యమని మీరు తీసుకోవచ్చు. అవి నిరూపించే విషయ పరిధిని పరిమితం చేస్తాయి; మనం మాట్లాడుతున్న వస్తువులకు పేర్లను కూడా పరిచయం చేస్తాయి. ఉదాహరణకు, మీ ప్రతిపాదన ``$!A$ అనేది $!B \lif !C$ రూపంలోని \tetoken{ఒక సూత్రం}{!!a{formula}} అనుకుందాం'' అని మొదలైతే, మీరు ఒక నిర్దిష్ట రకపు సూత్రాల (అంటే సోపాధికాల) గురించే నిరూపిస్తున్నారు; మీ నిరూపణలో చర్చించే $!B \lif !C$ ఏదైనా ఒక సోపాధికం అని అర్థం.

\end{document}
